\documentclass[11pt,a4paper]{article}

\usepackage[utf8]{inputenc}
\usepackage[T1]{fontenc}
\usepackage{lmodern}
\usepackage[french]{babel}
\usepackage[margin=1.9cm]{geometry}
\usepackage{microtype}
\usepackage{booktabs}
\usepackage{array}
\usepackage[table]{xcolor}
\usepackage{enumitem}
\usepackage{eurosym}
\usepackage[most]{tcolorbox}

\definecolor{navy}{HTML}{1F3864}
\definecolor{bluemid}{HTML}{2E5496}
\definecolor{lightblue}{HTML}{D6E0F0}
\definecolor{accent}{HTML}{C0491F}
\definecolor{softgreen}{HTML}{E2EFDA}
\definecolor{softamber}{HTML}{FCE4D6}
\definecolor{softred}{HTML}{F8D7DA}

\newtcolorbox{cle}[1][]{colback=lightblue!40,colframe=navy,boxrule=0.6pt,
  arc=2pt,left=8pt,right=8pt,top=6pt,bottom=6pt,#1}
\newtcolorbox{alerte}[1][]{colback=softamber!55,colframe=accent,boxrule=0.6pt,
  arc=2pt,left=8pt,right=8pt,top=6pt,bottom=6pt,#1}

\newcommand{\fait}{\cellcolor{softgreen}\textbf{fait}}
\newcommand{\partiel}{\cellcolor{softamber}partiel}
\newcommand{\trou}{\cellcolor{softred}\textbf{trou}}
\newcommand{\afaire}{à faire}

\setlength{\parskip}{3pt}

\title{\vspace{-1.4cm}\color{navy}\bfseries Feuille de route du mémoire\\[3pt]
  \large Objectif Prix SCOR : ce qui est fait, ce qui est fait mais invisible,
  et les trois trous}
\author{Kélian Kaddouri}
\date{État au 17 août 2026}

\begin{document}
\maketitle
\vspace{-0.8cm}

\begin{cle}
\textbf{Où en est le mémoire, en une phrase.} Le modèle est construit, calibré, gelé et
vérifié ; la rédaction couvre dix-neuf chapitres et cent cinquante-neuf pages. Ce qui reste
n'est donc plus de \emph{produire} mais de \textbf{rendre visible} ce qui est déjà mesuré, et de
fermer trois trous précis. La version de juillet de ce document est périmée : sa phase~3,
l'élicitation, présentée alors comme le principal apport restant, est \textbf{abandonnée} depuis
le 12~août, et la décision est documentée en annexe avec son chiffrage.
\end{cle}

\section*{1. L'état, mesuré plutôt que déclaré}

\begin{center}\footnotesize
\renewcommand{\arraystretch}{1.25}
\begin{tabular}{@{}p{5.4cm} r p{7.6cm}@{}}
\toprule
\textbf{Grandeur} & \textbf{Valeur} & \textbf{Comment elle se relit} \\
\midrule
Corps du mémoire & 117 p. & première page d'annexe moins une, dans \texttt{main.toc} \\
Document complet & 159 p. & compte des pages du PDF, jamais un index \\
Chapitres rédigés & 19 & dont 6 annexes \\
Scripts de calcul & 97 & \texttt{vasicek\_lab/*/*.py} \\
Sorties versionnées & 93 & \texttt{sorties\_verif/NN.txt} \\
Scripts cités par le mémoire & 74 & le reste est du travail non publié, voir la section 3 \\
Nombres publiés sous contrôle & 1\,865 & harnais, tous chapitres \\
Confirmation & 98,2 \% & 1\,831 confirmés, 34 dépouillés un par un \\
\textbf{Couverture} & \textbf{100 \%} & 0 nombre hors contrôle non déclaré, et c'est le
chiffre qui passe en premier \\
Contrôles du document & 0 / 0 / 0 / 0 & référence indéfinie, Overfull \verb|\vbox|,
annotation hors page, page tournée \\
\bottomrule
\end{tabular}
\end{center}

\section*{2. Ce qu'un jury de prix récompense, et où l'on se situe}

\begin{center}\footnotesize
\renewcommand{\arraystretch}{1.3}
\begin{tabular}{@{}p{4.3cm} c p{9.4cm}@{}}
\toprule
\textbf{Critère} & \textbf{État} & \textbf{Constat} \\
\midrule
Originalité maîtrisée & \fait &
La dépendance \emph{à l'ordre} entre piliers, absente de l'AMDEC comme des matrices de risque,
traduite en capital. Le mécanisme est neuf et il est borné, pas survendu. \\

Honnêteté méthodologique & \fait &
\textbf{C'est l'atout numéro un.} La non-transitivité qui donnait son titre au mémoire a été
réfutée par son propre auteur ; le chiffre central a été requalifié en borne supérieure ; la
borne inférieure de validité de la méthode est publiée ; les grandeurs simulées sortent avec
leur bruit. Peu de mémoires font cela. \\

Rigueur statistique & \fait &
Identification partielle ancrée (Manski, Tamer), valeur de Shapley ancrée depuis le 17~août
(Shapley, Owen, Song et al., Castro et al., Iooss et Prieur, Decupère), allocation d'Euler
ancrée (Tasche). \\

Pertinence réglementaire & \fait &
DORA, et la donnée manquante muée en spécification de registre. Directement opposable. \\

Reproductibilité & \partiel &
93 sorties versionnées, parité vérifiée entre deux machines, harnais de vérification de chaque
nombre publié. \textbf{Et le mémoire ne le dit nulle part}, voir la section 3. \\

Fiabilité inter-juges & \trou &
Le chapitre d'identifiabilité annonce une mesure d'accord entre deux codeurs indépendants. Le
second codage n'a pas eu lieu, le script~54 attend son CSV. \\

Ancrage sur entités réelles & \partiel &
Quatre bilans SFCR exploités, mais deux SCR sur quatre sont déduits d'un taux de couverture et
n'ont pas été relus dans les rapports. \\
\bottomrule
\end{tabular}
\end{center}

\section*{3. Les trois trous, par ordre de ce qu'ils coûtent}

\begin{alerte}
\textbf{Trou n\textsuperscript{o}1, et de loin le plus rentable : du travail fait et invisible.}
Vingt sorties versionnées ne sont citées par aucun chapitre, dont \textbf{les sept plus récentes}
(79, 80, 81, 83, 85, 86, 87). Les cinq notions empruntées au préprint de cascade climatique ont
été \emph{mesurées} et leurs résultats ne figurent dans aucun chapitre : ni la concavité du
capital en durée de non-conformité, ni l'équivalence observationnelle entre Hawkes et cascade,
ni le plafond lu comme amortisseur. Elles vivent dans les scripts et dans un deck, pas dans le
mémoire.

\textbf{Et le harnais lui-même n'est décrit nulle part.} Le mot n'apparaît qu'une fois dans le
corps, en passant, comme si le lecteur savait déjà de quoi il s'agit. Un dispositif qui vérifie
1\,865 nombres publiés contre les sorties des scripts qui les produisent, avec une couverture de
100~\%, est un \emph{résultat de méthode} et pas une coulisse. Aucun jury ne peut le créditer
s'il ne le lit pas.
\end{alerte}

\begin{alerte}
\textbf{Trou n\textsuperscript{o}2 : le second codage en aveugle.} C'est le seul endroit où le
mémoire annonce une mesure qui n'existe pas encore. Le kit est prêt, dix récits, une heure de
travail, et le formulaire est en état d'être envoyé depuis le 17~août. Tant qu'il manque, la
direction de $W$ repose sur un codeur unique, ce qui est précisément l'objection qu'un jury
formulera en premier sur la partie la plus originale du mémoire.
\end{alerte}

\begin{alerte}
\textbf{Trou n\textsuperscript{o}3 : le chiffre de tête n'est pas défini.} Trois valeurs
circulent pour \og l'écart entre conforme et non conforme\fg{} : $-53\,\%$ (script~25, à $g$ maintenu
au niveau non conforme), un facteur $3{,}34$ (script~43, quatre canaux) et $-64\,\%$
(script~20, sous OpRisk). Elles ne mesurent pas la même chose parce qu'elles ne déplacent pas
les mêmes paramètres. À une question de soutenance sur le résultat principal, il faut
\textbf{une} réponse et sa définition. Ne pas remplacer un de ces chiffres par un autre sans
avoir tranché : ce serait échanger un chiffre mal défini contre un autre.
\end{alerte}

\section*{4. La feuille de route}

\begin{center}\footnotesize
\renewcommand{\arraystretch}{1.3}
\begin{tabular}{@{}p{0.4cm} p{3.9cm} p{8.8cm} p{1.5cm}@{}}
\toprule
\textbf{Ph.} & \textbf{Objectif} & \textbf{Livrables} & \textbf{Statut} \\
\midrule
0 & Cadrage et thèse & Pivot vers la cascade dirigée ; architecture en cinq parties ;
plan Institut des Actuaires. & \fait \\

1 & Ancrage empirique & Fréquence (08b), sévérité POT gelée (07), accumulation et cas d'usage
(08e--08g), rejet argumenté du Hawkes contre les noyaux de la littérature (08h). & \fait \\

2 & Cascade et capital & Moteur \texttt{euro\_cascade\_model} ; SCR par état ; les quatre canaux
et leur interaction super-additive ; table complète des seize configurations (43, 68). & \fait \\

3 & Identifiabilité & Frontière explicite : la donnée fixe la co-occurrence, pas la direction.
Identification partielle sur les $2^{10}$ sommets. Biais de narration mesuré (64), séquences
ordonnées explorées et refermées avec critère de réouverture (75). & \fait \\

4 & Gel et vérification & Calibration gelée au 7~août. Harnais de vérification, couverture
portée de 74,7 à 100~\%. 93 sorties versionnées. Parité deux machines. & \fait \\

5 & Élicitation & \textbf{Abandonnée} le 12~août. Le protocole préparé et non exécuté devient
une annexe, avec le chiffrage de la raison : un panel mal calibré déplacerait le SCR de
777~M\euro{} sans qu'aucun signal ne l'indique. & \fait \\

6 & Robustesse et limites & Trois étages d'incertitude ; six postures avec leur bruit ; table à
deux colonnes des écarts involontaires et des choix prudents ; additivité des coûts bornée par
un exposant. & \fait \\

7 & Réponses au tutorat & Six demandes d'Hugo traitées (82, 83, 84) ; retours de Caroline
appliqués ; note de quatre pages. & \fait \\

8 & \textbf{Publier le travail mesuré} & Porter au mémoire les cinq notions (79--81, 85, 86), le
désaccord assumé avec le préprint sur le classement des leviers, et \textbf{une section
décrivant le dispositif de vérification}. & \trou \\

9 & \textbf{Fermer la fiabilité inter-juges} & Envoyer le formulaire, obtenir un second codage,
lancer le script~54, publier le coefficient d'accord. & \trou \\

10 & \textbf{Trancher le chiffre de tête} & Choisir la définition publiée de l'écart entre états
et l'imposer partout, decks compris. & \trou \\

11 & Consolidation finale & Vérifier les quatre jeux SFCR contre les rapports ; versionner la
sortie du script~35 ; choisir la posture reportée avec Caroline ; registre de sous-traitance. &
\partiel \\

12 & Soutenance & Narratif, anticipation des questions, banque de réponses (six entrées
aujourd'hui), relecture anti-jargon. & \afaire \\
\bottomrule
\end{tabular}
\end{center}

\section*{5. Chemin critique}

\textbf{Ce qui ne dépend que de moi, par ordre de rendement.} La phase~8 est la plus rentable du
tableau : le travail est fait, mesuré et versionné, il ne reste qu'à l'écrire, et elle
transforme sept scripts muets en autant de résultats lisibles. La section sur le dispositif de
vérification est le meilleur rapport valeur sur effort de tout le projet, parce qu'elle rend
créditable tout le reste. La phase~10 est une décision, pas un calcul, et elle se prend en une
heure.

\textbf{Ce qui dépend des autres, donc à lancer en premier.} Le second codage (phase~9) est le
seul élément à délai externe. Tout le reste peut attendre, celui-là non : il faut envoyer,
relancer, et prévoir qu'une partie des sollicités ne répondra pas. Le formulaire est prêt depuis
le 17~août, le mail aux contacts de l'Institut et le message à la squad cyber sont rédigés.

\textbf{Ce qui est clos et ne doit pas se rouvrir.} Le gel de la calibration, la convention de
normalisation de $W$ par l'émission maximale, le statut de citation de Hackmageddon, le rejet du
Hawkes, la non-transitivité, l'élicitation, la convention qui sépare la VaR de la charge annuelle
du quantile de sévérité d'un sinistre, et le détecteur de contradictions entre scripts, instruit
puis écarté pour une raison.

\vspace{0.25cm}
\begin{cle}
\textbf{Le piège de juillet a été évité, il en reste deux.} L'avertissement de la version
précédente était de ne pas survendre $W$ comme calibré : il a été tenu, et au-delà, puisque la
frontière d'identifiabilité est devenue une contribution à part entière.

\textbf{Premier piège restant : diluer le registre d'honnêteté.} Ce qui tient ce mémoire est
qu'il publie ses propres limites. Toute reformulation qui rendrait une réserve moins visible le
dégrade, même si elle le fait paraître plus assuré. Retirer un $\pm$ pour faire net, effacer que
la posture la plus prudente est la moins précise, taire un désaccord avec une source citée par
ailleurs comme convergente : trois façons de perdre ce qui distingue le travail.

\textbf{Second piège, et c'est celui du moment : laisser du travail mesuré hors du document.}
Un résultat qui n'est pas dans le mémoire n'existe pas pour le jury. Sept scripts et un
dispositif de vérification complet sont aujourd'hui dans ce cas.
\end{cle}

\vspace{0.2cm}
\noindent\footnotesize\textbf{Note de calendrier.} Aucune date de remise ni de soutenance ne
figure dans le dépôt : ce document est donc ordonné par \emph{rendement}, pas par échéance. Les
phases 8 à 10 sont à traiter avant la phase 12, et la phase 9 est à lancer aujourd'hui puisque
son délai n'est pas le mien.

\vspace{0.15cm}
\noindent\footnotesize\textbf{Source des chiffres du \S1.} Comptes relevés le 17 août 2026 sur
le PDF compilé, sur \texttt{main.toc}, sur l'arborescence de \texttt{vasicek\_lab} et sur la
sortie de \texttt{verif\_tous\_chapitres.ps1}. Aucun n'est recopié d'une version antérieure de
ce document, la précédente ayant précisément péri de cela.

\end{document}
